% Keep references off the technical pages. When content occupies four pages,
% the next page is references only, as required by EMBC 2027.
\clearpage
\IEEEtriggeratref{6}
\begin{thebibliography}{11}
\raggedright

\bibitem{schalk2007}
G.~Schalk \emph{et al.}, ``Decoding two-dimensional movement trajectories using electrocorticographic signals in humans,'' \emph{J. Neural Eng.}, vol.~4, no.~3, pp.~264--275, 2007, doi:~\href{https://doi.org/10.1088/1741-2560/4/3/012}{10.1088/1741-2560/4/3/012}.

\bibitem{miller2019}
K.~J. Miller, ``A library of human electrocorticographic data and analyses,'' \emph{Nature Human Behaviour}, vol.~3, pp.~1225--1235, 2019, doi:~\href{https://doi.org/10.1038/s41562-019-0678-3}{10.1038/s41562-019-0678-3}.

\bibitem{competition2008}
K.~J. Miller and G.~Schalk, ``Prediction of finger flexion: 4th brain-computer interface data competition,'' dataset description, Jun.~2008. [Online]. Available: \url{https://www.bbci.de/competition/iv/desc_4.pdf}.

\bibitem{competitionresults}
BCI Competition IV organizers, ``BCI Competition IV: Final results,'' Dataset~4. Accessed: Oct.~9, 2026. [Online]. Available: \url{https://www.bbci.de/competition/iv/results/index.html}.

\bibitem{fingerflex}
V.~Lomtev, A.~Kovalev, and A.~Timchenko, ``FingerFlex: Inferring finger trajectories from ECoG signals,'' \emph{arXiv preprint arXiv:2211.01960}, version~2, 2023. [Online]. Available: \url{https://arxiv.org/abs/2211.01960v2}.

\bibitem{dtcnet}
F.~Wang, Z.~Luo, W.~Lv, and X.~Zhu, ``DTCNet: Finger flexion decoding with three-dimensional ECoG data,'' \emph{Front. Comput. Neurosci.}, vol.~19, Art.~no.~1627819, 2025, doi:~\href{https://doi.org/10.3389/fncom.2025.1627819}{10.3389/fncom.2025.1627819}.

\bibitem{switching2016}
A.~Elgharabawy and M.~Abdel Wahed, ``Decoding of finger movement using kinematic model classification and regression model switching,'' in \emph{Proc. 8th Cairo Int. Biomed. Eng. Conf. (CIBEC)}, 2016, pp.~84--89, doi:~\href{https://doi.org/10.1109/CIBEC.2016.7836126}{10.1109/CIBEC.2016.7836126}.

\bibitem{bc4d4}
G.~Jangir, N.~Joshi, and G.~Purohit, ``Harnessing the synergy of statistics and deep learning for BCI competition 4 dataset 4: A novel approach,'' \emph{Brain Informatics}, vol.~12, Art.~no.~5, 2025, doi:~\href{https://doi.org/10.1186/s40708-025-00250-5}{10.1186/s40708-025-00250-5}.

\bibitem{hilofusenet}
Q.~Sun, E.~Calvo Merino, B.~Van Dyck, Y.~Yang, J.~He, and M.~M. Van Hulle, ``Spectro-temporal fusion of high-gamma and low-frequency ECoG signals for intracranial finger movement decoding,'' author manuscript, May~26, 2026. [Online]. Available: \url{https://github.com/hisunjiang/HiLoFuseNet/blob/bee83100953c85332e9ad24020345d8dad79a963/manuscript.pdf}.

\bibitem{corteg}
L.~Yang, Q.~Sun, B.~Van Dyck, E.~Calvo Merino, and M.~M. Van Hulle, ``CORTEG: Foundation models enable cross-modality representation transfer from scalp to intracranial brain recordings,'' \emph{arXiv preprint arXiv:2605.10337}, version~1, 2026. [Online]. Available: \url{https://arxiv.org/abs/2605.10337v1}.

\bibitem{trace2026}
Y.~Yu, ``TRACE: Transformer for Regularized and Accurate Cortical ECoG Motor Decoding,'' author-released benchmark repository, 2026, commit~\texttt{5c7e47a}. Accessed: Oct.~9, 2026. [Online]. Available: \href{https://github.com/epyifany/TRACE/tree/5c7e47a469e266b472dfec82cf4adc379b25c23a}{\texttt{github.com/epyifany/TRACE}}.

\end{thebibliography}
